\chapter{The Philosophy of Representation}
\label{ch:philosophy-of-representation}

Have you ever thought about words like "gold" or "silver"? Or why, when we see a tree, we call it a "tree," even though the tree itself bears no resemblance whatsoever to the letters and sound of that word?\\
Why can the number "5" refer equally to five objects, to 5 people, or to five years?\\
Why, when we say "the February Revolution," do we not have to reconstruct every detail and event of it in our minds?\\
When we think about these things, it seems our brain and our language have found a strange way of coping with the world: instead of retaining the exact details of objects and events, we work with symbols and names for things.\\
At the end of the previous chapter we asked a key question: "How can a word, a number, or a sign take the place of something else in this way?" The goal of this chapter is to answer exactly that question. In fact, this chapter deals with one of the most basic aspects of human nature: naming, and its role in remembering information.
\newpage
\section{Why Are Humans So Fond of Naming?}
\label{sec:why-naming}
Naming and symbol-making have accompanied humans for as long as we can trace, and they still do, from drawing shapes on cave walls to, today, drawing in Photoshop. But why? Why are humans so keen on putting a name on everything?\\
The simplest answer we can arrive at is "laziness." We have to admit that the human race is far lazier than it thinks it is, and however lazy we are, our brains are lazier still. To make this clearer, imagine you are the manager of a large warehouse that receives tens of thousands of long letters every day, and your job is to find and deliver any requested letter as fast as possible. What do you think the best approach would be? Before reading on, take a moment to think it through and find the best solution for yourself, and only then keep going.\\
Our brain is always facing exactly this kind of situation. Every day it has to store, manage, and retrieve thousands upon thousands of things in short-term and long-term memory. But how?

\subsection{The Brain's Warehouse Manager: The Hippocampus}
\label{sub:hippocampus-warehouse-manager}
One of the brain's key components for managing memory is a small structure called the hippocampus, which plays the role of warehouse manager in our analogy. The Cleveland Clinic explains that the hippocampus is the part of the brain responsible for memory and learning, and that this small structure helps store material, both short-term and long-term, and build awareness of one's surroundings\autocite{clevelandclinic-hippocampus}. In other words, the hippocampus plays an important role in converting short-term memories into long-term ones, and in organizing and retrieving them.\\
Working alongside this "warehouse manager" is a network responsible for retrieving and managing those memories: episodic memory, the ability to encode and retrieve one's own daily personal experiences, which is supported by the medial temporal lobe (MTL)\autocite{roy2018hippocampalindex}. Put simply, the hippocampus acts like a manager who organizes the data and "letters" of memory, while the episodic memory network acts like a postal system tasked with collecting and retrieving those memories.\newpage
An important point about the brain's memory-retrieval system is that the brain does not preserve every detail of an event in full. Instead, it mostly keeps an index, a marker, or a general symbol of the event. For example, if you once had cake smeared on your face at a birthday party and it made you excited or upset, that moment has likely been stored as the general symbol of that birthday in your brain and in the minds of those around you. In other words, our memory is not like a complete archive of scenes; it is a collection of memory traces that gets presented to us\autocite{mitnews2017circuit}.\\
That said, our aim here is not to go into finer neuroscientific detail; it was simply to get a general sense of how the brain's storage function works, to show that memory is essential to everyday life and involves a range of specific abilities that, at their core, allow information to be stored and retrieved across time spans ranging from seconds to days to years. It is also now well established that memory takes several distinct forms\autocite{frankland2019retrieval}.\\
And in practice, when we try to recall something like our "age of 9," the entirety of that period does not appear before our eyes. Rather, only a few standout points, a fear, an excitement, a particular moment from childhood, come to mind, and we tell ourselves, "ah, that belongs to when I was 9." Based on these facts about the brain, we arrive at something called "representation" and "mental representation."
\section{Mental Representation}
\label{sec:mental-representation}
Before defining representation itself, we can define mental representation, a notion that goes back at least to Aristotle. It takes as its starting point common-sense mental states such as thoughts, beliefs, desires, perceptions, and imaginings. Such states are said to have "intentionality": they are about things, or refer to things, and can be evaluated along dimensions such as consistency, truth, appropriateness, and accuracy. For instance, the thought that cousins are not related is inconsistent; the belief that Elvis is dead is true; the desire to eat the moon is inappropriate; the visual experience of a ripe strawberry as exactly red is accurate; and picturing George Washington with braided hair is inaccurate\autocite{sep-mental-representation}.
\newpage Mental representation also treats mental processes such as thinking, reasoning, and imagining as sequences of intentional mental states. For example, imagining the moon rising over a mountain involves, among other things, entertaining a set of mental images of the moon and of a mountain\autocite{sep-mental-representation}.\\
From this angle, we realize that even in recalling memories we are dealing with mental representations, not with the past itself. For instance, when we say "tree," are we really seeing the tree itself, or the mental image that word has constructed in our mind? Or when we talk about "the February Revolution," do all the historical events and images actually replay in our minds, or is it simply a name we have chosen to represent that set of events?\\
Many questions could be asked here, and we arrive at something rather strange: in order to manage our world, we do not necessarily carry the world itself, with all its details, along with us. For a tree we have a name, for a quantity a number, for a historical event a phrase, and even for our own memories, what returns to mind is not necessarily the event itself.\\
And if you look closely, all of these put one thing in the place of another. This is exactly what we have been talking about from the start: "representation," of which mental representation is a very precise example.

\newpage

\section{Two Classical Views}
\label{sec:two-classical-views}
Now it's time to go back to the root of the matter and answer a deeper, more important question:

"When we say `tree,' do we see the tree itself, or do we merely see something that represents it?"

This question, in all its modern formulations, has very old roots. To see one of the earliest serious statements of the problem, we can turn to two philosophers bound by a relationship of teacher and student, who nonetheless took very different paths when it came to imitation and representation: Plato and Aristotle.

Let's start with Plato.

In Book X of the \textit{Republic}, while discussing painting and imitation, Plato draws a distinction between a thing itself and the appearance, or image, that gets made of it. In a dialogue with Glaucon, Socrates offers a simple but crucial example: if you take a mirror in your hand and turn it around, you can see in it the sun, the earth, yourself, animals, plants, and even objects like a bed. But has the mirror actually made these things?

Glaucon's answer is clear: no; the mirror only brings forth their "appearance," not their reality and truth.

Plato then places the painter in this same category. A painter can paint a picture of a bed, but what he has produced is not the bed itself; it is only an appearance of it. This is why, when Socrates asks Glaucon what the painter actually makes, the answer points toward "appearance," not reality itself.

A few lines later Plato sharpens the question further: does painting, in every case, imitate reality as it is, or appearance as it appears?

Plato's own answer is the second.

This distinction matters for our discussion, because here we run into something that is not the thing itself, yet points to the thing and can place an appearance of it before us. The picture of the bed is not the bed itself; yet without the bed we also cannot explain what this picture is "a picture of."

From this point Plato extends his critique to imitative art in general. For him the issue isn't simply that the image doesn't match the real object; the issue is that representation can confront us with the appearance of a thing while the thing itself, and the truth about it, remains something else entirely.

So on this view, there is a gap between "the thing" and "what we see of the thing" --- a gap that should not be overlooked.

But the story doesn't end here.

Aristotle, Plato's student, takes a different path when facing "imitation" and what appears to us of a thing. In the \textit{Poetics}, unlike Plato's critical stance toward imitative art, he stresses that imitation is a natural part of human nature, and that from childhood human beings learn their very first things through it. Aristotle also explains that people take pleasure in seeing representations, because in observing them they can learn and recognize what each image points to --- for instance, "this is that person." This discussion appears in the \textit{Poetics}, 1448b4--19.

Here, representation is no longer simply something that opens a gap between us and reality. It can be a means by which we recognize something and learn something. This doesn't mean Aristotle takes every representation to be identical with the truth; rather, his point is that our relationship with whatever represents a thing can itself be a source of knowledge.

But Aristotle takes the question one step deeper still.

In \textit{On the Soul} (\textit{De Anima}), while examining the difference between perception, imagination, and thought, he introduces the concept of \textit{phantasia}. In Book III, Chapter 3, Aristotle states explicitly that imagination is not the same as sense perception or discursive thought. He even illustrates this with the example that a person can call up an "image" in the mind, the same way certain mnemonic techniques make use of mental images. He then concludes this discussion by treating \textit{phantasia} as a motion that arises out of the actualization of a sense faculty. In this way, something that arose within sense experience can go on to play a role in mental activity even in a process separate from perception itself.

But in Chapter 4 Aristotle goes even more fundamental. He asks how the faculty with which we think can come to know anything at all. His answer is that the thinking part of the soul must be able to receive the "form" of a thing without itself becoming that thing --- that is, it must be potentially identical with its object of thought, without ever actually becoming that object. In this same discussion Aristotle describes the soul as "the place of forms," with the qualification that these forms are not, at first, actual within it, but only potential.

\subsection{Plato: The Cave and Imperfect Representations}
\label{sub:plato-cave}

But let's look more deeply at these two teachers' views on representation.

One of the most famous images Plato left behind in the whole history of philosophy is the "Allegory of the Cave," which he presents in Book VII of the \textit{Republic}. Imagine a group of people chained since childhood in the depths of a cave, in such a way that they can only stare at the wall in front of them. Behind them a fire burns, and between the fire and the prisoners, other people carry objects and statues whose shadows fall on the wall. The prisoners, having never seen anything else, take these shadows to be the only reality there is\autocite{plato-republic}.

The important point for our discussion is this: the cave's prisoners are doing exactly what --- as we saw in earlier sections --- our own minds do all the time as well: they put a representation (the shadow) in the place of the thing itself, without even noticing this substitution. The difference is that Plato sees this condition not as a useful tool, but as a prison.

This suspicion toward representation comes through even more clearly in Book X of the \textit{Republic}. There Plato turns to the example of a bed: the bed that exists in the Form, made by god; the carpenter who, looking at that Form, makes one particular physical bed; and the painter who, without knowing anything about carpentry or even about the bed itself, paints only the bed's appearance from one particular angle onto canvas. In Plato's view, the painter's work is "three removes from the truth": once from the eternal Form of the bed to the physical bed, and again from the physical bed to the painted image\autocite{plato-republic}. For Plato, then, representation --- whether a painting or a shadow on the cave wall --- is always something diminished and incomplete: a copy of a copy, growing more deceptive the further it strays from its original source.

\subsection{Aristotle: The Useful Internal Model}
\label{sub:aristotle-model}

Aristotle, Plato's student, parts ways with his teacher precisely at this point. He rejected the idea of Forms existing in some separate realm above the world of the senses; in his view, the form of a thing is not located somewhere else, but lies within the thing itself --- a position later called hylomorphism (the theory of form and matter), which Aristotle set explicitly against Plato's theory of Forms\autocite{sep-form-matter}. For Aristotle, when we say "this is a bed," the form of "bedness" has to be found in the wood and nails themselves, not in some separate heavenly realm as Plato imagined.

This metaphysical difference has a direct consequence for Aristotle's view of representation. In \textit{On the Soul} (\textit{De Anima}) he speaks of a faculty of the soul called \textit{phantasia}, which builds mental images and makes them available to the intellect; without these images, he holds, thought is not possible at all. Aristotle even has a famous line for this: "the soul never thinks without an image"\autocite{aristotle-deanima}. In other words, for Aristotle mental images are not a prison or a departure from the truth, but a necessary tool.

Aristotle spells this idea out more precisely in another treatise, \textit{On Interpretation} (\textit{De Interpretatione}). In its opening lines he writes that spoken words are symbols (\textit{symbola}) for affections of the soul, and written marks are in turn symbols for those spoken words; but unlike languages, which differ from people to people, those inner affections are the same for everyone, because they are all likenesses (\textit{homoiomata}) of the very same external things\autocite{aristotle-deinterpretatione}. This short passage can be read as the first clear statement of a chain this book will return to again and again: the word represents the mental state, and the mental state represents the thing itself.

So if we set these two classical views side by side: in representation, Plato sees a distancing from the truth; Aristotle sees, in that very same representation, the only possible route to thought. What's striking is that this second view is exactly what we already rediscovered empirically at the start of this chapter, when we talked about the hippocampus and memory management: our brain builds representations of the world in place of the world itself, not because it is lazy or defective, but because it has no other option.

\newpage
\section{The Modern Philosophical Encounter with Representation}
\label{sec:modern-encounter}

Several centuries later, the very same question that had emerged in ancient Greek philosophy, about the relation between a thing and its image or appearance, took on a new shape. This time the question was no longer just how closely an image resembles the thing it represents; it was pushed back into perception and cognition themselves:

When we see something, what exactly are we encountering? The thing itself, or something the mind holds in its place?

Among the philosophers who reshaped this path in the most fundamental way, René Descartes and Immanuel Kant hold a special place. Descartes, with his method of doubt, shook our naive trust in the senses and showed that what seems obvious in everyday experience is not necessarily known as simply as we assume. Kant, roughly a century and a half later, took a further step: rather than merely asking whether our knowledge corresponds to things, he asked what conditions make it possible for us to experience and know things at all.

\subsection{Descartes: Doubting Sensory Representations}
\label{sub:descartes-doubt}

In the \textit{Meditations on First Philosophy} (1641), René Descartes launches an ambitious project: setting aside every belief that admits even the smallest room for doubt, so that, in the end, something wholly certain might remain. The first casualty of this doubt is the senses. Descartes reminds us that the senses have deceived us before, and so don't deserve complete trust; he even entertains the possibility that we might be dreaming, or that a deceiving demon might be fabricating all our perceptions. Under this radical doubt, every sensory representation that once seemed self-evident suddenly becomes suspect\autocite{descartes-meditations}.

But Descartes carries this doubt forward with a very concrete example known as the "wax argument." Consider a fresh piece of wax: it smells of honey, tastes of nectar, has a definite color and shape, and is cold and hard. Now bring it near a fire: the smell disappears, the taste fades, the shape changes, and the wax turns warm, soft, even liquid --- none of its original sensory qualities remain. And yet we immediately say this is "the same wax." Descartes asks: what, then, shows us this sameness of the wax, when the senses have given us nothing but a shifting, contradictory set of qualities? His answer: not the senses, and not even the imagination, but only the intellect --- the mind's pure judgment --- that recognizes the wax as one persisting thing\autocite{descartes-meditations}. So even our simplest, most everyday representation of an object --- the representation of a "persisting" object --- is an achievement of the mind, not a transparent gift from the senses.

Descartes extends this insight, in the Third Meditation, into a fuller theory of representation. He examines every "idea," or mental conception, from two angles: on one side, an idea is simply a mental event that exists --- he calls this the idea's "formal reality"; on the other side, an idea represents something, and he calls that represented content the idea's "objective reality"\autocite{sep-descartes-ideas}. My idea of a stone and my idea of God, for instance, are identical in formal reality --- both are simply states of my mind --- but wildly different in objective reality, that is, in what they represent. This may be the first time in the history of philosophy that "mental representation" is formulated this clearly, as an independent, analyzable subject in its own right.

\subsection{Kant: The Mind Imposes a Framework}
\label{sub:kant-framework}

If Descartes showed that even the simplest sensory representation requires the mind's judgment, Immanuel Kant, roughly a century and a half later, took a far bigger step: he showed that no experience whatsoever exists for us prior to the mind's involvement.

In the \textit{Critique of Pure Reason} (first edition 1781, revised second edition 1787), in the preface to the second edition, Kant compares his change of approach to a Copernican revolution. Just as Copernicus, instead of assuming the heavens revolve around a stationary observer, set the observer in motion, Kant proposes that instead of assuming our cognition must conform to objects, we assume that objects must conform to the conditions of our cognition\autocite{kant-critique}. This is a complete inversion: before Kant, everyone thought of the mind as a mirror reflecting the outside world; Kant says no --- it is the mind that gives shape to the world.

Specifically, Kant argues that space and time are not properties of things in themselves at all, but merely "pure forms of sensible intuition" --- frameworks that our mind imposes on raw sense data prior to any experience. In addition, the understanding possesses a set of a priori concepts called the "categories" (such as causality and substance), which lend coherence and meaning to this already organized sensory data. At one point Kant even writes explicitly that we must regard all appearances as "mere representations," not as things that exist in themselves\autocite{sep-kant-idealism}. What we experience --- the "appearance" --- is always a product of this mental framing; and the "thing in itself," what exists independent of this framework, remains forever beyond our reach.

The consequence of this view matters a great deal for our discussion. For Plato, and even for Descartes, representation remained a problem to be solved: is our representation correct or false? Does it correspond to a mind-independent reality or not? But for Kant this question no longer carries its earlier meaning, because, on his account, there simply is no experience of "reality" that isn't already mediated by representation. Representation, on this view, is no longer a barrier standing between us and the world; it is the very condition that makes having a world for us possible at all.

\section{Four Criteria for Representation}
\label{sec:four-criteria-representation}

But having come this far, one question remains unanswered: when we say something represents something else, what exactly makes us call it a "representation"? Does anything that points to something else count as a representation? For instance, a piece of paper on a desk can remind you of an appointment, but is it itself a representation?\\
Philosophy and cognitive science have wrestled with exactly this question for years. One attempt to clarify the matter, particularly within neuroscience, proposes four criteria for the concept of representation\autocite{baker2021representation}. These four criteria are not meant to be a final definition agreed upon by every philosopher; rather, they can help us understand what exactly we mean when we talk about a representation.\newpage

\subsection{Connection to the Real World}
\label{sub:criterion-real-world-connection}
First, a representation must have some connection, in the real world, to the thing it represents. This means we must be able to find some relationship between the "representer" and "the thing represented." If no relationship exists between the two at all, it becomes difficult to say one represents the other.
\subsection{A Causal Role in Behavior}
\label{sub:criterion-causal-role}
Second, a representation must be able to play a causal role in behavior or in some process. That is, a representation must not merely be correlated with or accompany something; it must be able to play a role in explaining what an organism or a system does. For example, if the brain represents something as "danger," that representation must play a role in explaining behavior such as fleeing.
\subsection{Exclusivity of the Causal Role}
\label{sub:criterion-exclusivity}
Third, a representation must uniquely occupy its causal role: the absence of the representation should block the behavior it is capable of producing. That is, if several different parts or processes could play exactly the same role in producing a behavior, it becomes harder to say that one of them specifically is the "representation" we mean. For this reason, in neuroscience, the mere fact that a pattern of activity correlates with a feature of the environment is not, by itself, sufficient; the specific role that pattern plays in the mechanism generating the behavior must also be clear.
\newpage
\subsection{Teleological Justification}
\label{sub:criterion-teleological}
Fourth, a representation must have a teleological justification that identifies the cases in which the representation is incorrect. This criterion tells us that if something is truly a representation, we must be able to explain what that representation is supposed to do, and under what conditions we can say it has gotten something "wrong." A map, for instance, can show a route correctly or incorrectly; a sentence can be true or false. So, in order to count something as a representation, we need a criterion that lets us distinguish a successful representation from a mistaken one.
\newpage
\section{One Thing in Place of Another}
\label{sec:something-in-place-of-another}
In this chapter we saw how our brain and language use representation to work with the world. Instead of holding on to every detail of a thing, we build something that stands "in place of" that thing. A word or a number, for example, is not the object itself, but a stand-in for it. Likewise, a map or a photograph is not the world itself, but marks its place.\\
You have probably heard something like this as a child, from a school principal or a teacher: "You are a mirror of your parents." Well, you are not literally a mirror of your parents, but you may reflect their behaviors, habits, manner of speech, or even some of their traits in yourself. Even in everyday language, we constantly use something that is not, itself, the other thing, yet somehow connects us to it.\\
In fact, what all these examples share is exactly this: one thing standing in the place of another. Now that we have understood this much, let me leave you with a question as an exercise: is representation only an abstract concept in our minds, or can it also exist in a physical system? And if it is to exist in the real world, how does it move from being an abstract concept to something physical and measurable?

\newpage
\section{The Layered Nature of Representation}
\label{sec:layered-representation}

By this point --- whether we side with Plato, who saw representation as a deviation from truth, or with Aristotle, who saw it as a necessary tool; whether with Descartes, who showed that even our simplest perception requires the mind's judgment, or with Kant, who cast the mind not as a passive receiver but as the active constructor of experience --- in all of these views there's one shared point that perhaps hasn't yet been stated outright: representation is never a flat, single-layer relation.

What does this mean? When Kant says the mind first pours raw sensory data into the molds of space and time and then organizes it further with categories like causality, he is, in effect, describing a layered structure: a raw, sensory layer at the bottom, a layer of perceptual forms above it, and a layer of the understanding's concepts above that. And when Aristotle says a word symbolizes a state of the soul, and a state of the soul is a likeness of a thing, we again run into a chain of layers: thing, mental state, word --- each layer representing something that is itself one layer below it, in the chain of representation.

This idea, that representation happens in "layers," has also found a precise formulation in contemporary cognitive science. David Marr, the neuroscientist and vision scientist, proposes in his influential book \textit{Vision} that any information-processing system --- whether the human brain or a computer --- can be analyzed at three distinct levels: the computational level (what problem is this system solving, and why), the algorithmic-representational level (with what representations and what processes does it solve that problem), and the implementational level (in what specific physical hardware --- neurons, say, or electronic circuits --- do these representations and processes actually get realized)\autocite{marr1982}. The key point for our discussion is this: the abstract representation described at the algorithmic level never stays suspended in midair; it must, in the end, be realized at that third level, in something physical and measurable.

If we return to the data-information-knowledge-wisdom (DIKW) pyramid we saw in the first chapter (Section~\ref{sec:levels-of-data}), we find this same layered structure again: each level is built on the one before it and represents that level at a higher degree of abstraction. And this is exactly what happens in our own memory as well: our memory of "being nine years old" is not a representation of that whole period, but a representation of a few standout memory traces from it; and those memory traces, at a still lower level, are themselves nothing more than patterns of neural activity.

So here is the subtle point this layered view brings out: when we ask, "what does this representation point to?", the answer is not always some raw, unmediated object out in the world. Very often, what one layer of representation points to is simply the layer of representation below it, not a world free of representation altogether. The word points to the concept, the concept to the perception, and the perception to the neural signal --- and how far does this chain go?

This is where the chapter's closing question takes on a sharper form. It is no longer enough to ask "is representation abstract or physical?", because we now know that representation happens across multiple layers, some of which look entirely abstract (a number, a concept, a theory). The real question is this: however long this downward chain of layers runs, where must it finally come to rest? And once it does, what does that final point look like?

The answer, as Marr's third level already hints, can be nothing other than the physical world: a voltage in a neuron, ink on paper, a pattern of magnetic field on a disk. In other words, every abstract representation, no matter how many layers of abstraction it has passed through, must in the end be carried by a physical signal.

And it is exactly here that the present chapter must come to a stop, giving way to the question the next chapter must answer:

\begin{quote}
\itshape
How does a representation, which up to this point we have understood merely as an abstract stand-in for something else, ultimately come to be carried by a physical signal?
\end{quote}